\section{Method Overview}
The implemented system combines signal processing, machine learning, and state-based supervision in MATLAB. The software baseline focuses on public IMU datasets, reproducible training and evaluation scripts, and replay-based inspection of the assist-gating behavior. Mechanical design and hardware actuation remain outside the submitted repository scope.

\section{Intention Detection Algorithm}
The repository performs classification on overlapping windows of IMU data with 100 samples per window and a 50-sample step. The main baseline is an \textbf{RBF SVM} classifier operating on handcrafted window-level features.

Per IMU, five features are computed: mean and variance of accelerometer magnitude, mean and variance of gyroscope magnitude, and dominant frequency of the vertical-acceleration spectrum. HuGaDB contributes up to six IMUs per window through \texttt{FeaturesFromImuStack}, producing a \textbf{30-dimensional} feature vector. USC-HAD contributes a single smartphone IMU and is zero-padded through \texttt{LocomotionFeatureVector} so both datasets share the same input dimension.

For deployment-oriented evaluation, the binary task maps activity labels to \textbf{active movement} versus \textbf{inactive}. The HuGaDB-centered path is the main one used for replay and assist gating. For exploratory analysis, the repository also includes dataset-native multiclass recognition using ECOC models over the native activity labels of each dataset.

An optional \textbf{binary LSTM} baseline processes the same windows as sequences rather than feature rows. Sliding windows are encoded as feature-by-time matrices with \texttt{ImuWindowToSequenceMatrix.m}, and training sequences are built with \texttt{PrepareTrainingDataSequences.m}. The LSTM is used as a standalone sequence model, not as a fused stage on top of the SVM.

\section{Multi-Sensor Data Processing Pipeline}
The HuGaDB path preserves multi-IMU structure and session metadata. The loader distinguishes \texttt{single\_activity} and \texttt{multi\_activity\_sequence} sessions, applies official gyro-corruption metadata, and filters invalid data through a shared quality gate before training or replay. In the cleaned GitHub v1 corpus, the \texttt{single\_activity} family still contains several concrete session variants such as \texttt{walking}, \texttt{running}, \texttt{sitting}, \texttt{standing}, \texttt{bicycling}, \texttt{sitting\_in\_car}, and \texttt{down\_by\_elevator}; the loader keeps those files available but canonicalizes them into the controlled \texttt{single\_activity} protocol family for filtering, while \texttt{various} recordings are treated as \texttt{multi\_activity\_sequence} sessions for the streaming pipeline.

Kalman-based processing is implemented through MATLAB's \texttt{imufilter} and wrapped in \texttt{FusionKalman.m}. In the submitted baseline, this fusion block is used for replay-time segment-pitch visualization rather than mandatory classifier preprocessing. This keeps the learning pipeline simple and reproducible while still allowing kinematic traces to be shown alongside the control output.

\section{Model Training and Evaluation}
Binary SVM training uses \texttt{fitcsvm} with standardized features. Multiclass recognition uses \texttt{fitcecoc}. Five-fold out-of-fold evaluation is used for the reported SVM results, with predictions generated through cross-validation utilities such as \texttt{crossval} and \texttt{kfoldPredict}.

The main scripts are \texttt{TrainSvmBinary}, \texttt{TrainSvmMulticlass}, \texttt{EvaluateSvmConfusion}, \texttt{EvaluateMulticlassConfusion}, and \texttt{RunSvmDatasetAblation}. The default deployment-oriented checkpoint is \texttt{Binary\_SVM\_Model\_hugadb\_only.mat}.

Optional sequence baselines are trained with the Deep Learning Toolbox. \texttt{TrainLstmBinary} and \texttt{TrainLstmMulticlass} build two-layer LSTM classifiers on sequence windows and evaluate them on separate stratified holdout splits. Their metrics are reported separately from the SVM baselines because their evaluation protocol differs from the five-fold OOF setup used for the main reported results.

\section{System Integration and Control}
Replay-time control logic is implemented as a binary assist supervisor rather than a torque controller. Window-level predictions are passed to \texttt{RealtimeFsm.m}, which applies asymmetric hysteresis: three consecutive active predictions are required to engage assistance, and five consecutive inactive predictions are required to release it. This reduces rapid toggling when predictions fluctuate near transition boundaries.

\texttt{RunExoskeletonPipeline.m} executes the binary SVM replay path on held-out HuGaDB sessions. \texttt{RunExoskeletonPipelineLstm.m} mirrors the same replay structure with the binary LSTM baseline. \texttt{RunExoskeletonPipelineMulticlass.m} provides an exploratory multiclass replay path through \texttt{RealtimeFsmFromActivityClass.m}, and \texttt{RunExoskeletonPipelineMulticlassLstm.m} runs the same multiclass control loop with the sequence LSTM classifier. In all cases, the implemented output is a stable binary assist command for simulation and visualization.

\section{List of MATLAB Libraries and Toolboxes}
\subsection*{Components in active use}
\begin{table}[htbp]
  \centering
  \caption{MATLAB components used in the implemented baseline.}
  \begin{tabular}{@{}ll@{}}
    \toprule
    \textbf{Component} & \textbf{Role} \\
    \midrule
    MATLAB & Core scripting and visualization \\
    Signal Processing Toolbox & FFT and signal operations \\
    Statistics and Machine Learning Toolbox & SVM / ECOC training and cross-validation \\
    Sensor Fusion and Tracking Toolbox & \texttt{imufilter} fusion \\
    Deep Learning Toolbox & Optional binary and multiclass LSTM models \\
    \bottomrule
  \end{tabular}
\end{table}
